\section{The Transit Integration Layer}

\subsection{Overview}

The transit integration layer adds \$2 billion to the I2.0 program --- 0.8\% of total capital --- and unlocks 9 T1/T1 diamond hubs and approximately 50 T1/T2 regional stops for intercity bus service. This section summarizes the transit case; the full analysis appears in the companion F.1 paper \citep{ROUTE_E1}.

The transit layer consists of passenger-oriented infrastructure at freight-optimized locations:
\begin{itemize}
  \item \textbf{Bus platforms with weather protection}: covered waiting areas adjacent to existing rest area facilities, sized for 4--8 intercity coaches
  \item \textbf{Passenger parking}: 200--400 spaces at T1/T1 hubs; 50--150 spaces at T1/T2 regional stops
  \item \textbf{Ticketing and information kiosks}: digital boarding pass readers and real-time departure boards integrated with carrier scheduling systems
  \item \textbf{ADA-accessible restrooms}: incremental upgrade to existing rest area facilities
  \item \textbf{Wi-Fi and device charging}: standard amenity for modern intercity travel
  \item \textbf{Safety lighting upgrade}: perimeter lighting consistent with passenger facility standards
\end{itemize}

\subsection{Why Near-Zero Marginal Cost}

The reason the transit layer is cost-effective at \$2B is that the hard infrastructure it depends on is already funded by the freight program. The diamond interchange upgrade (Component 2, \$4.5B) provides grade-separated access roads, utility connections, stormwater management, and site security. The rest area improvements funded by the managed lane program (Component 1) provide overnight parking, restrooms, and food service. The transit platform increment --- the physical addition of covered platforms, passenger parking, and ticketing --- costs \$1.8--2.5 million per T1/T1 hub, compared to \$20--50 million for a standalone urban bus terminal \citep{BTS_intermodal2023}.

At standalone terminal cost, the same transit coverage would cost approximately \$400 million for the 9 T1/T1 hubs and \$2.5 billion for the 50 T1/T2 regional stops --- a total of \$2.9 billion for the hubs alone, against the \$59 million hub-increment cost in the I2.0 program. The savings from colocation with freight infrastructure approach 50:1 on a per-hub basis.

\subsection{City Pairs Newly Connected}

The F.1 analysis identifies 40 or more origin-destination pairs where both cities are within 30 miles of a T1 hub and currently have no direct intercity bus service with fewer than one transfer. The Greyhound bankruptcy of 2020 eliminated service on approximately 1,200 route-segments in the United States, concentrated in the Southeast, Gulf Coast, and Midwest \citep{Amtrak2023}. FlixBus US expansion has restored service on some high-density corridors but has not penetrated rural or small-metro markets where load factors are uncertain.

The T1 hub network makes bus operator entry viable on routes that were previously uneconomical because the physical infrastructure (a safe, accessible boarding point with parking) did not exist. Bus operators cannot serve a highway rest area without passenger facilities; they can serve a T1 hub with a platform and shelter. The hub network effectively lowers the capital threshold for bus operator market entry on underserved corridors.

\subsection{Equity Dimension}

The F.1 paper shows that T1/T1 hub locations align disproportionately with communities scoring high on the C3 dimension (Economic Opportunity Access) --- corridors in which the 30-mile buffer population has below-average GDP per capita and above-average transit-dependent household rates \citep{ACS_B08201_2022}. In 7 of the 9 T1/T1 hub locations, the C3 score of the surrounding area exceeds 5.0 (where 5.0 represents the national median). The hubs in Gary/Chicago (south side), San Antonio, and the rural T1/T2 stops along I-69 and I-29S show the strongest alignment between highway location and transit need.

This alignment is not coincidental. The rubric design embeds C3 (Economic Opportunity Access) as a scoring dimension precisely because corridors serving economically disadvantaged areas have the strongest claim to public investment. The result is that the rubric-driven site selection for diamond hubs, driven primarily by freight intensity and connectivity arguments, also produces the right answer for transit equity.

\subsection{Transit Integration as Mandatory Design Standard}

The F.1 paper concludes --- and this synthesis endorses --- that transit integration should be a mandatory design standard for all T1/T1 diamond hub construction, not an optional addition. The argument is straightforward: if a diamond hub is designed and constructed without transit-compatible platform geometry, passenger-accessible utilities, or ADA-compliant access paths, the cost of retrofitting those features post-construction is 3--5 times higher than building them in from the start. Road geometry must be designed for bus turning radii. Utility conduits must be routed for passenger facility connections. Parking must be designed with pedestrian connections to the platform area.

A diamond hub constructed to freight-only standards cannot be efficiently upgraded for transit service. The incremental cost to include transit design from the outset is below 1\% of hub construction cost. The incremental cost to retrofit a freight-only hub for transit service is 15--25\% of the original hub cost. The recommendation follows directly: transit integration is an investment protection measure for the \$4.5B diamond upgrade program, not an add-on cost.
